\documentclass[addpoints, 11pt]{exam}
\usepackage[a4paper, left=20mm, right=20mm, top=26mm, bottom=30mm, headsep=6mm, footskip=12mm, headheight=15pt]{geometry}
\usepackage[T1]{fontenc}
\usepackage[utf8]{inputenc}
\usepackage{lmodern}
\usepackage{textcomp}
\usepackage[french]{babel}
\usepackage{xcolor}
\usepackage{listings}
\usepackage{graphicx}
\usepackage{array}
\usepackage{amsmath}
\usepackage{eurosym}
\usepackage{needspace}
\usepackage[most]{tcolorbox}
\usepackage{tikz}
\usetikzlibrary{positioning,arrows.meta}


% --- Code Python : même style que le sujet, en police à chasse fixe ---
\definecolor{codegreen}{rgb}{0,0.6,0}
\definecolor{codegray}{rgb}{0.5,0.5,0.5}
\definecolor{codepurple}{rgb}{0.58,0,0.82}
\definecolor{backcolour}{rgb}{0.95,0.95,0.92}
\lstdefinestyle{mystyle}{
    backgroundcolor=\color{backcolour},
    commentstyle=\color{codegreen},
    keywordstyle=\color{magenta},
    numberstyle=\tiny\color{codegray},
    stringstyle=\color{codepurple},
    basicstyle=\footnotesize\ttfamily,
    columns=fullflexible,
    breakatwhitespace=false,
    breaklines=true,
    postbreak=\mbox{\textcolor{codegray}{$\hookrightarrow$}\space},
    captionpos=b,
    keepspaces=true,
    numbers=left,
    numbersep=5pt,
    xleftmargin=1.5em,
    showspaces=false,
    showstringspaces=false,
    showtabs=false,
    tabsize=4,
    upquote=true
}
\lstset{style=mystyle, language=Python, extendedchars=true,
  literate={à}{{\`a}}1 {â}{{\^a}}1 {ä}{{\"a}}1 {ç}{{\c{c}}}1
           {é}{{\'e}}1 {è}{{\`e}}1 {ê}{{\^e}}1 {ë}{{\"e}}1
           {î}{{\^i}}1 {ï}{{\"i}}1 {ô}{{\^o}}1 {ö}{{\"o}}1
           {ù}{{\`u}}1 {û}{{\^u}}1 {ü}{{\"u}}1 {œ}{{\oe}}2
           {É}{{\'E}}1 {È}{{\`E}}1 {À}{{\`A}}1 {Ç}{{\c{C}}}1
           {°}{{\textdegree}}1 {€}{{\euro}}1 {…}{{\dots}}1
           {«}{{\guillemotleft}}1 {»}{{\guillemotright}}1 {·}{{\textperiodcentered}}1}
% Terminal (commandes Linux) : sans couleur de mots-clés
\lstdefinestyle{terminal}{style=mystyle, language={}, numbers=none, xleftmargin=0.5em}
% Sortie d'un programme (ce qui s'affiche à l'écran)
\lstdefinestyle{sortie}{style=mystyle, language={}, numbers=none, backgroundcolor=\color{white}, frame=single, rulecolor=\color{black!60}, xleftmargin=0.5em}

% --- Boîtes du corrigé (lisibles en noir et blanc) ---
\newtcolorbox{reponse}[1][Réponse attendue]{enhanced, breakable, colback=white, colframe=black, boxrule=1.2pt, arc=0pt,
  fonttitle=\bfseries, coltitle=white, colbacktitle=black, title=#1, left=2mm, right=2mm, top=1mm, bottom=1mm, before skip=6pt, after skip=6pt}
\newtcolorbox{erreur}[1][Erreur fréquente]{enhanced, breakable, colback=black!6, colframe=black!55, boxrule=0.6pt, arc=2pt,
  borderline west={3pt}{0pt}{black!70}, fonttitle=\bfseries, coltitle=black, colbacktitle=black!15, title=#1, left=2mm, right=2mm, top=1mm, bottom=1mm, before skip=6pt, after skip=6pt}
\newtcolorbox{astuce}[1][À retenir]{enhanced, breakable, colback=white, colframe=black!55, boxrule=0.6pt, arc=2pt,
  fonttitle=\bfseries, coltitle=black, colbacktitle=black!8, title=#1, left=2mm, right=2mm, top=1mm, bottom=1mm, before skip=6pt, after skip=6pt}
\newtcolorbox{remarque}[1][Remarque sur l'énoncé]{enhanced, breakable, colback=white, colframe=black!70, boxrule=0.6pt, arc=0pt,
  borderline={0.4pt}{1.2pt}{black!70, dashed}, fonttitle=\bfseries, coltitle=black, colbacktitle=white, title=#1, left=2mm, right=2mm, top=1mm, bottom=1mm, before skip=6pt, after skip=6pt}
\newcommand{\enonce}[1]{\textit{#1}\par\smallskip}
\newcommand{\methode}{\par\smallskip\noindent\textbf{Méthode.}\ }
\newcommand{\explication}{\par\smallskip\noindent\textbf{Explication.}\ }
\newcommand{\code}[1]{\texttt{#1}}
% Coupure de ligne possible après un tiret bas (noms de variables longs)
\DeclareRobustCommand{\_}{\textunderscore\allowbreak}
% Un code court ne se coupe pas en bas de page
\AddToHook{env/lstlisting/before}{\needspace{10\baselineskip}}
\setlength{\parindent}{0pt}
\setlength{\parskip}{3pt}
\renewcommand{\arraystretch}{1.25}

\title{Corrigé --- CC1 1re NSI - Binaire, variables et conditions (2023)}
\author{Babash - NSI Tani Malandi et Johanne Bakalara}
\date{Septembre 2023, durée : 1 h}

\begin{document}
\pagestyle{headandfoot}
\runningheader{}{}{\footnotesize Corrigé --- Binaire (2023)}
\firstpageheader{}{}{}
\runningfooter{}{\footnotesize Babash - NSI Tani Malandi et Johanne Bakalara \textperiodcentered{} CC BY-SA 4.0 \textperiodcentered{} page \thepage}{}
\firstpagefooter{}{\footnotesize Babash - NSI Tani Malandi et Johanne Bakalara \textperiodcentered{} CC BY-SA 4.0 \textperiodcentered{} page \thepage}{}

\maketitle
\thispagestyle{headandfoot}

\begin{astuce}[Comment lire ce corrigé]
Pour chaque question : la question est rappelée en italique, puis la \textbf{réponse attendue} est encadrée en noir, puis l'\textbf{explication} montre la méthode, étape par étape.
Les encadrés gris « Erreur fréquente » signalent les pièges les plus courants. Essaie d'abord de refaire la question seul·e, puis compare ta démarche avec celle du corrigé.
\end{astuce}

\section*{Barème}
\begin{center}
  \gradetable[h]
\end{center}

\begin{questions}

% =====================================================================
\section{Conversion octets-bits}

\question[1] \enonce{Sam a un téléphone de 80 Go. Une musique pèse environ 240 Mbits. Combien de musiques Sam peut-il stocker ?}

\begin{reponse}
$80 \text{ Go} = 80 \times 10^9 \times 8 = 640 \times 10^9$ bits et $240 \text{ Mbits} = 240 \times 10^6$ bits.\\
$\dfrac{640 \times 10^9}{240 \times 10^6} \approx 2\,666{,}7$ : Sam peut stocker \textbf{2\,666 musiques complètes}.
\end{reponse}

\explication
On cherche « combien de fois » une musique rentre dans la mémoire : c'est une \textbf{division}. Mais on ne divise que deux quantités \textbf{de même unité}. Or la mémoire est en \textbf{octets} (Go) et la musique en \textbf{bits} (Mbits).

\methode
\begin{enumerate}
  \item Mettre la mémoire en bits. 1 Go $= 10^9$ octets et 1 octet $=8$ bits :\\
  $80 \text{ Go} = 80 \times 10^9 \text{ octets} = 80 \times 10^9 \times 8 \text{ bits} = 640 \times 10^9 \text{ bits}$.
  \item Mettre la musique en bits. 1 Mbit $= 10^6$ bits : $240 \text{ Mbits} = 240 \times 10^6 \text{ bits}$.
  \item Diviser :
  $\dfrac{640 \times 10^9}{240 \times 10^6} = \dfrac{640}{240} \times 10^{3} = \dfrac{8}{3} \times 1\,000 \approx 2\,666{,}7$.
  \item Interpréter : on ne stocke pas 0,7 musique. On garde la \textbf{partie entière} : 2\,666 musiques complètes.
\end{enumerate}
\textit{Autre chemin, aussi juste :} mettre la musique en Mo. $240 \text{ Mbits} = 240 \div 8 = 30$ Mo et $80 \text{ Go} = 80\,000$ Mo ; $80\,000 \div 30 \approx 2\,666{,}7$.

\begin{erreur}
Oublier le facteur 8 entre bits et octets : on trouve $80\,000 \div 240 \approx 333$ musiques, 8 fois trop peu. Regarde toujours si l'unité est un \textbf{b} (bit) ou un \textbf{o} (octet).
\end{erreur}

\question[1] \enonce{Combien de valeurs différentes peut-on stocker sur 16 bits ?}

\begin{reponse}
$2^{16} = \mathbf{65\,536}$ valeurs différentes (par exemple les entiers de 0 à 65\,535).
\end{reponse}

\explication
Chaque bit a \textbf{2} valeurs possibles : 0 ou 1. Avec 1 bit : 2 valeurs. Avec 2 bits : $2 \times 2 = 4$ valeurs (00, 01, 10, 11). Chaque bit ajouté \textbf{double} le nombre de possibilités. Avec 16 bits : $2 \times 2 \times \dots \times 2 = 2^{16} = 65\,536$.

\begin{erreur}
Répondre 65\,535. C'est la plus grande valeur, pas le nombre de valeurs : de 0 à 65\,535, il y a 65\,536 nombres, car on compte aussi le 0.
\end{erreur}

% =====================================================================
\section{Décimal, binaire et hexadécimal}

\question[1] \enonce{Convertir $63_{10}$ en binaire en détaillant les calculs. Expliquer pourquoi ce résultat était facile à calculer.}

\begin{reponse}
$63_{(10)} = 32 + 16 + 8 + 4 + 2 + 1 = \mathbf{11\,1111_{(2)}}$.\\
C'était facile car $63 = 64 - 1 = 2^6 - 1$ : un nombre de la forme $2^n - 1$ s'écrit avec $n$ chiffres 1.
\end{reponse}

\methode\ (les soustractions successives)
\begin{enumerate}
  \item Écrire les puissances de 2 : 128, 64, 32, 16, 8, 4, 2, 1.
  \item Prendre la plus grande puissance qui \textbf{ne dépasse pas} le nombre, la soustraire, et écrire 1 dans sa colonne.
  \item Recommencer avec le reste. Écrire 0 dans les colonnes non utilisées.
\end{enumerate}
\begin{center}
\begin{tabular}{|l|c|}
\hline
Calcul & Reste \\ \hline
$63 - 32 = 31$ & 31 \\
$31 - 16 = 15$ & 15 \\
$15 - 8 = 7$ & 7 \\
$7 - 4 = 3$ & 3 \\
$3 - 2 = 1$ & 1 \\
$1 - 1 = 0$ & 0 : fini \\ \hline
\end{tabular}
\qquad
\begin{tabular}{l|cccccccc}
 & $2^7$ & $2^6$ & $2^5$ & $2^4$ & $2^3$ & $2^2$ & $2^1$ & $2^0$\\
 & 128 & 64 & 32 & 16 & 8 & 4 & 2 & 1\\ \hline
$63 =$ & 0 & 0 & 1 & 1 & 1 & 1 & 1 & 1\\
\end{tabular}
\end{center}
On lit le résultat de gauche à droite, sans les 0 du début : $11\,1111_{(2)}$.

\textbf{Pourquoi c'était prévisible.} $64 = 2^6$ s'écrit $100\,0000_{(2)}$ (un 1 suivi de six 0). En retirant 1, on obtient le nombre juste en dessous : six chiffres 1, comme $999 = 1000 - 1$ en base 10.

\textit{Autre méthode :} les divisions successives par 2. $63 = 2 \times 31 + \mathbf{1}$ ; $31 = 2 \times 15 + \mathbf{1}$ ; $15 = 2 \times 7 + \mathbf{1}$ ; $7 = 2 \times 3 + \mathbf{1}$ ; $3 = 2 \times 1 + \mathbf{1}$ ; $1 = 2 \times 0 + \mathbf{1}$. On lit les restes \textbf{de bas en haut} : $111111$.

\question[1] \enonce{Convertir $113_{16}$ en décimal.}

\begin{reponse}
$113_{(16)} = 1 \times 256 + 1 \times 16 + 3 \times 1 = \mathbf{275_{(10)}}$.
\end{reponse}

\methode
\begin{enumerate}
  \item En base 16, chaque position vaut une puissance de 16 : $16^0 = 1$, $16^1 = 16$, $16^2 = 256$.
  \item Placer les chiffres de droite à gauche sous ces puissances.
  \item Multiplier chaque chiffre par sa puissance, puis additionner.
\end{enumerate}
\begin{center}
\begin{tabular}{l|ccc}
 & $16^2$ & $16^1$ & $16^0$\\
 & 256 & 16 & 1\\ \hline
$113_{(16)} =$ & 1 & 1 & 3\\
\end{tabular}
\end{center}
$1 \times 256 + 1 \times 16 + 3 \times 1 = 256 + 16 + 3 = 275$.

\begin{erreur}
Lire « 113 » comme cent treize. L'indice $_{16}$ dit que ce nombre est écrit en base 16 : le premier 1 vaut 256, pas 100.
\end{erreur}

\question[2] \enonce{Décoder le message binaire en ASCII (conseil : passer par la base 16).}

\begin{reponse}
Le message est : \textbf{« Vous marquez 2 pts ! »}
\end{reponse}

\explication
En ASCII, chaque caractère est codé par \textbf{un octet} (8 bits). Le tableau du sujet donne le caractère à partir de son code \textbf{hexadécimal} : le premier chiffre hexadécimal donne la ligne, le second la colonne.

\methode\ (pour chaque octet)
\begin{enumerate}
  \item Couper l'octet en deux paquets de 4 bits.
  \item Convertir chaque paquet en un chiffre hexadécimal (de 0 à F), avec les poids 8, 4, 2, 1.
  \item Lire le caractère dans le tableau : ligne $=$ premier chiffre, colonne $=$ second chiffre.
\end{enumerate}
Exemple avec le premier octet : $0101\,0110$. $0101 = 4 + 1 = 5$ et $0110 = 4 + 2 = 6$. Le code est $56_{(16)}$ : ligne 5, colonne 6, c'est \textbf{V}.

\begin{center}
\small
\begin{tabular}{|c|c|c||c|c|c|}
\hline
Octet & Hexa & Car. & Octet & Hexa & Car. \\ \hline
0101 0110 & 56 & V & 0111 0001 & 71 & q \\
0110 1111 & 6F & o & 0111 0101 & 75 & u \\
0111 0101 & 75 & u & 0110 0101 & 65 & e \\
0111 0011 & 73 & s & 0111 1010 & 7A & z \\
0010 0000 & 20 & espace & 0010 0000 & 20 & espace \\
0110 1101 & 6D & m & 0011 0010 & 32 & 2 \\
0110 0001 & 61 & a & 0010 0000 & 20 & espace \\
0111 0010 & 72 & r & 0111 0000 & 70 & p \\
 & & & 0111 0100 & 74 & t \\
 & & & 0111 0011 & 73 & s \\
 & & & 0010 0000 & 20 & espace \\
 & & & 0010 0001 & 21 & ! \\ \hline
\end{tabular}
\end{center}
Le tableau se lit colonne de gauche, puis colonne de droite.

\begin{erreur}
Confondre le chiffre 2 (code $32_{(16)}$) avec le nombre 2. Le code $32_{(16)}$ est le \textbf{caractère} « 2 ». Autre piège : oublier les espaces (code $20_{(16)}$), qui sont des caractères comme les autres.
\end{erreur}

% =====================================================================
\section{Complément à 2}

\question[3] \enonce{Question de cours et réflexion.}
\begin{parts}
\part \enonce{À quoi sert le complément à 2 ?}
\begin{reponse}
Le complément à 2 sert à \textbf{représenter les entiers relatifs} (positifs \textbf{et négatifs}) en binaire, dans la mémoire de l'ordinateur.
\end{reponse}
Ses avantages : les nombres positifs gardent leur écriture habituelle, le zéro a une seule écriture, et l'addition se fait avec le même algorithme que pour les entiers positifs. Le processeur n'a donc pas besoin d'un circuit spécial pour les négatifs.

\part \enonce{Dire si chaque affirmation est vraie ou fausse.}
\begin{reponse}
\begin{enumerate}
  \item Les positifs gardent leur représentation classique : \textbf{Vrai}.
  \item Un nombre plus son opposé vaut 0 : \textbf{Vrai}.
  \item Le zéro a deux représentations : \textbf{Faux}.
  \item L'algorithme d'addition reste le même : \textbf{Vrai}.
\end{enumerate}
\end{reponse}
\explication
\begin{enumerate}
  \item Sur 8 bits, $5$ s'écrit $0000\,0101$, comme en binaire « normal ».
  \item $5 = 0000\,0101$ et $-5 = 1111\,1011$. En posant l'addition, on obtient $1\,0000\,0000$ : 9 bits. Le 9\textsuperscript{e} bit (la retenue) ne tient pas sur 8 bits ; il est perdu. Il reste $0000\,0000 = 0$.
  \item Sur 8 bits, le seul zéro est $0000\,0000$. Le nombre $1000\,0000$ vaut $-128$, pas « $-0$ ». (C'est la méthode « bit de signe + valeur absolue » qui a deux zéros : c'est justement un défaut que le complément à 2 corrige.)
  \item Voir le 2 : on additionne bit à bit avec les retenues, comme d'habitude.
\end{enumerate}

\part \enonce{Par quel symbole commencent les nombres positifs et les nombres négatifs ?}
\begin{reponse}
Les positifs (et zéro) commencent par \textbf{0} ; les négatifs commencent par \textbf{1}.
\end{reponse}
Le premier bit, à gauche, est appelé \textbf{bit de poids fort}. En complément à 2, il indique le signe.
\end{parts}

\question[1] \enonce{Quelle est la valeur décimale de $1001\,0010_{(2)}$, représenté en complément à 2 (sur 8 bits) ?}

\begin{reponse}
$1001\,0010_{(2)} = -128 + 16 + 2 = \mathbf{-110}$.
\end{reponse}

\methode\ (le tableau du complément à 2)
\begin{enumerate}
  \item Sur 8 bits, le bit de gauche vaut $-2^7 = -128$ au lieu de $+128$. Les autres gardent leur poids habituel.
  \item Additionner les poids des bits à 1.
\end{enumerate}
\begin{center}
\begin{tabular}{cccccccc}
$-2^7$ & $2^6$ & $2^5$ & $2^4$ & $2^3$ & $2^2$ & $2^1$ & $2^0$\\
$-128$ & 64 & 32 & 16 & 8 & 4 & 2 & 1\\ \hline
1 & 0 & 0 & 1 & 0 & 0 & 1 & 0\\
\end{tabular}
\end{center}
$-128 + 16 + 2 = -110$.

\textit{Vérification par l'autre méthode} (inverser les bits, puis ajouter 1) : le nombre commence par 1, il est négatif. On inverse : $0110\,1101$ ; on ajoute 1 : $0110\,1110 = 64 + 32 + 8 + 4 + 2 = 110$. Le nombre vaut donc $-110$.

\begin{erreur}
Calculer $128 + 16 + 2 = 146$ : c'est la valeur si l'octet code un entier \textbf{positif} (non signé). L'énoncé dit « complément à 2 » : le bit de gauche compte pour $-128$.
\end{erreur}

\end{questions}

\end{document}
